\section{Proofs for the boundary results}\label{app:lim}

This appendix proves \cref{thm:lim:ambient,thm:lim:posslp,prop:lim:flow}.

\subsection{Proof of \texorpdfstring{\cref{thm:lim:ambient}}{the ambient-count theorem}}\label{app:lim:ambient}

\paragraph{One block.}
Let $n=m^2$, which is even. Choose $u\in\Q^n$ with $n/2$ entries $1/m$ and
$n/2$ entries $-1/m$, such that $u_{n-1}=u_n=1/m$, and let
$B_+,B_-\subseteq\{1,\ldots,n-2\}$ be the index sets of the remaining positive
and negative entries, so $\abs{B_+}=n/2-2$ and $\abs{B_-}=n/2$. Then
$\norm u=1$. Put $q=e_{n-1}-e_n$, which is orthogonal to $u$, and
$c=\frac34$, $\varepsilon=m^{-3}$,
\[
 A_1=\alpha\bigl[c(uq^{\mathsf T}+qu^{\mathsf T})+qq^{\mathsf T}\bigr],
\]
and let $P_1\subset\R^n$ be the polytope defined by
\[
 x_i\ge0\ (i\le n-2),\qquad\sum_{i\le n-2}x_i=4m,\qquad
 x_{n-1}+x_n=0,\qquad\abs{x_{n-1}-x_n}\le\varepsilon .
\]
The instance of the theorem consists of $k$ copies:
$A=\diag(A_1,\ldots,A_1)$, $T=\diag(u^{\mathsf T},\ldots,u^{\mathsf T})$, and
$P=P_1^k$. Its data have encoding length polynomial in $N$.

\emph{Part \ref{it:lim:amb-instance}.} In the orthonormal basis
$(u,q/\sqrt2)$ of their common range, $A_1$ and $A_1+\alpha uu^{\mathsf T}$
act as
\[
 \alpha\begin{pmatrix}0&c\sqrt2\\c\sqrt2&2\end{pmatrix}
 \qquad\text{and}\qquad
 \alpha\begin{pmatrix}1&c\sqrt2\\c\sqrt2&2\end{pmatrix},
\]
and both vanish on the orthogonal complement. The second matrix has positive
trace and determinant $\alpha^2(2-2c^2)=\frac78\alpha^2>0$, so
$A+\alpha T^{\mathsf T}T\succeq0$. The first has eigenvalues
$\alpha(1\pm\sqrt{1+2c^2})=\alpha(1\pm\sqrt{17/8})$, exactly one of them
negative, and $\nu=\alpha(\sqrt{17/8}-1)>\alpha/4$ because
$\sqrt{17/8}>\frac54$. The rows of $T$ have disjoint supports and unit norm,
so $TT^{\mathsf T}=I_k$. For $x\in P_1$ write $y=u^{\mathsf T}x$ and
$z=q^{\mathsf T}x$. Then $x_{n-1}=z/2$, $x_n=-z/2$, and
$y=m^{-1}(\sum_{B_+}x_i-\sum_{B_-}x_i)$; since the nonnegative $x_i$,
$i\le n-2$, have sum $4m$, the pair $(y,z)$ ranges exactly over
$[-4,4]\times[-\varepsilon,\varepsilon]$. Moreover
$\frac12x^{\mathsf T}A_1x=\alpha(cyz+z^2/2)$.

\paragraph{The auxiliary function of one block.}
Let $\gamma\in\R^n$, $\xi=u^{\mathsf T}\gamma$, $\zeta=\gamma-u\xi$, and
\[
 S=\frac m2\Bigl(\min_{B_+}\gamma_i-\min_{B_-}\gamma_i\Bigr),\qquad
 t=\frac{\gamma_{n-1}-\gamma_n}2 .
\]
We claim that, for a number $K$ depending on $\gamma$ only,
\begin{equation}\label{eq:lim:amb-block}
 V(a)=K+\varphi(a),\qquad
 \varphi(a)=\min_{\abs y\le4,\ \abs z\le\varepsilon}
 \Bigl[G(y,z)+\frac\alpha2\Bigl(a+\frac\xi\alpha-y\Bigr)^2\Bigr],
\end{equation}
with $G(y,z)=\alpha(cyz+z^2/2)+Sy+tz$,
where $V$ is the function \eqref{eq:lim:aux} for one block. Fix $(y,z)$ and
minimize $\zeta^{\mathsf T}x$ over the points of $P_1$ with these
projections. The last two coordinates are determined and contribute
$(\zeta_{n-1}-\zeta_n)z/2=tz$, because $u_{n-1}=u_n$. The
masses $\sum_{B_+}x_i=m(4+y)/2$ and $\sum_{B_-}x_i=m(4-y)/2$ are determined,
and a linear function on a scaled simplex is minimized at a vertex, so the
minimum is $\frac{m(4+y)}2\min_{B_+}\zeta_i+\frac{m(4-y)}2
\min_{B_-}\zeta_i$. Since $\zeta_i=\gamma_i-\xi/m$ on $B_+$ and
$\zeta_i=\gamma_i+\xi/m$ on $B_-$, this equals $K_0+(S-\xi)y$ with
$K_0$ independent of $(y,z)$. Hence
\[
 V(a)=K_0+\xi a+\min_{y,z}\Bigl[\alpha\Bigl(cyz+\frac{z^2}2\Bigr)
 +(S-\xi)y+tz+\frac\alpha2(a-y)^2\Bigr],
\]
and $\frac\alpha2(a-y)^2+\xi(a-y)=\frac\alpha2(a-y+\xi/\alpha)^2-\xi^2/(2\alpha)$
gives \eqref{eq:lim:amb-block} with $K=K_0-\xi^2/(2\alpha)$.

The objective inside \eqref{eq:lim:amb-block} is a quadratic in $(y,z)$ with
Hessian $\alpha\bigl(\begin{smallmatrix}1&c\\c&1\end{smallmatrix}\bigr)\succ0$.
It has a unique minimizer $(y^*(a),z^*(a))$ on the rectangle, and by
Danskin's theorem $\varphi$ is differentiable with
$\varphi'(a)=\alpha(a+\xi/\alpha-y^*(a))$.

\paragraph{A good event.}
Let $\mathcal E_1$ be the event that $\min_{B_+}\gamma_i\le-1+4/n$,
$\min_{B_-}\gamma_i\le-1+4/n$, and $\abs\xi\le2$. We show
$\Prob(\mathcal E_1)>\frac14$ under both laws. Let
$\pi=\Prob\{\gamma_i\le-1+4/n\}$. For the uniform law $\pi=2/n$. For
$U_{1,M}$ the atoms $-1+2j/(M-1)$ in this range are those with
$j\le2(M-1)/n$, at least $2(M-1)/n$ of them, so
$\pi\ge(2/n)(1-1/M)\ge19/(10n)$ because $M\ge32$. In both cases
$\pi\ge19/(10n)$. Using $1-\pi\le e^{-\pi}$, $\abs{B_+}\ge15n/32$ (as
$n\ge64$), and $\abs{B_-}=n/2$,
\[
 \Prob\Bigl\{\min_{B_+}\gamma_i\le-1+\tfrac4n\Bigr\}\ge1-e^{-57/64}>\tfrac7{12},
 \qquad
 \Prob\Bigl\{\min_{B_-}\gamma_i\le-1+\tfrac4n\Bigr\}\ge1-e^{-19/20}>\tfrac35,
\]
because $e^{57/64}>\frac{12}5$ and $e^{19/20}>\frac52$. These two events
depend on disjoint coordinates and are independent. Next,
$\operatorname{Var}\xi=\sum_iu_i^2\operatorname{Var}\gamma_i=\operatorname{Var}
\gamma_1$, which is $\frac13$ for the uniform law and $(M+1)/(3(M-1))\le\frac38$
for the grid law, so $\Prob\{\abs\xi>2\}\le\frac3{32}$ by Chebyshev's
inequality. Therefore $\Prob(\mathcal E_1)\ge\frac7{12}\cdot\frac35-\frac3{32}
=\frac{41}{160}>\frac14$. No independence between $\xi$ and the two minima is
used.

\paragraph{Flatness on the good event.}
On $\mathcal E_1$, both minima lie in $[-1,-1+4/n]$, so
$\abs S\le\frac m2\cdot\frac4n=\frac2m$; also $\abs t\le1$ and
$\abs\xi\le2$.

\emph{Small derivative.} Let $\abs a\le\frac32$. For every
$\abs z\le\varepsilon$, the unconstrained minimizer in $y$ of the objective in
\eqref{eq:lim:amb-block} is $a+\xi/\alpha-(S+\alpha cz)/\alpha$, whose
absolute value is at most $\frac32+2+\frac2m+\varepsilon<4$ because
$\alpha\ge1$. So $y^*(a)$ is this point, and
$\varphi'(a)=S+\alpha cz^*(a)$. Since $\alpha\le m$,
\begin{equation}\label{eq:lim:amb-flat}
 \abs{V'(a)}=\abs{\varphi'(a)}\le\frac2m+\frac34\,m\cdot m^{-3}<\frac3m
 \qquad(\abs a\le\tfrac32).
\end{equation}

\emph{Small gap.} Let $\abs a\le1$. The point $(y,z)=(a+\xi/\alpha,0)$ is
feasible, since $\abs y\le3$, so $\varphi(a)\le G(a+\xi/\alpha,0)\le3\abs S\le
6/m$. Every feasible $(y,z)$ has
$G(y,z)\ge-4\alpha c\varepsilon-4\abs S-\varepsilon\ge-3m^{-2}-8m^{-1}-m^{-3}$,
and $\min\varphi\ge\min G$. Hence
\begin{equation}\label{eq:lim:amb-gap}
 V(a)-\min V=\varphi(a)-\min\varphi<\frac{15}m\qquad(\abs a\le1).
\end{equation}

\paragraph{Products and counts.}
For the $k$-block instance, $A$, $T$, and $P$ split into blocks, so the
minimization in \eqref{eq:lim:aux} splits and
$V(a)=\sum_{l=1}^kV^{(l)}(a_l)$, where $V^{(l)}$ is the one-block
function of the $l$th block of noise coordinates. The block events
$\mathcal E_1^{(l)}$ are independent, so their intersection
$\mathcal E$ has probability greater than $4^{-k}$. The allowance is
$2E_h=\alpha kh^2/4$.

For (b)(i), let $8/\sqrt{\alpha m}\le h\le16/\sqrt{\alpha m}$, so
$2E_h\ge16k/m$. On $\mathcal E$, every $v\in\mathbb G_h\cap[-1,1]^k$ has
$V(v)-\min V=\sum_l(V^{(l)}(v_l)-\min V^{(l)})<15k/m<2E_h$ by
\eqref{eq:lim:amb-gap}. A grid of spacing $h$ has at least
$2/h-1\ge\sqrt{\alpha m}/8-1\ge\sqrt{\alpha m}/16$ nodes in $[-1,1]$, because
$\alpha m\ge256$. So the expected count is at least
$4^{-k}(\sqrt{\alpha m}/16)^k=(\sqrt{\alpha m}/64)^k$.

For (b)(ii), let $16/(\alpha m)\le h\le32/(\alpha m)$. Then $h\le\frac18$, so
$v_l\pm h\in[-\frac32,\frac32]$ for $v_l\in[-1,1]$. On $\mathcal E$,
\eqref{eq:lim:amb-flat} and the mean value theorem give, for every $l$ and
both signs,
\[
 V(v)-V(v\pm he_l)=V^{(l)}(v_l)-V^{(l)}(v_l\pm h)
 \le\frac{3h}m\le\frac{\alpha h^2}4\le2E_h ,
\]
because $h\ge12/(\alpha m)$. The grid has at least
$2/h-1\ge\alpha m/16-1\ge\alpha m/32$ nodes in $[-1,1]$, so the expected count
is at least $4^{-k}(\alpha m/32)^k=(\alpha m/128)^k$.

For (c), $\alpha=1$ and $m=\sqrt{N/k}$ turn the two bounds into
$64^{-k}(N/k)^{k/4}$ and $128^{-k}(N/k)^{k/2}$. The matrix $A$ has $O(N)$
nonzero entries of $O(\log N)$ bits and $P$ has $O(N)$ constraints, so even
a dense encoding has $I\le C_0N^2\log N$ for an absolute constant $C_0$. The
quantities recorded in $R$ do not change with $m$. If either expected count
were at most $f(k,R)I^c$, then fixing $k>8c$ and letting $m$ grow would
contradict $64^{-k}(N/k)^{k/4}\le f(k,R)(C_0N^2\log N)^c$.

Finally, a search grid of the low-rank route qualifies: it divides an
auxiliary box that contains $[-\frac32,\frac32]$ in every coordinate by
powers of two, and its meshes $w2^{-j}$ meet every interval $[t,2t]$ with
$t$ at most the box width $w$, in particular both ranges of $h$ in (b). The
ambient diameter of $P_1$ is at least $4\sqrt2\,m$, since $P_1$ contains
$4me_i$ for every $i\le n-2$; the theorem therefore does not contradict
bounds that grow with the ambient diameter.\qed

\subsection{Proof of \texorpdfstring{\cref{thm:lim:posslp}}{the PosSLP theorem}}\label{app:lim:posslp}

A straight-line program is a sequence of integers $A_1,\ldots,A_s$ in which
each $A_j$ is $0$, $1$, or $A_l\circ A_m$ with $l,m<j$ and
$\circ\in\{+,-,\times\}$; its output is $A=A_s$. The integers can have
exponentially many bits. The construction never computes them; it writes
only bounded local formulas.

\begin{lemma}[Bounded pair encoding]\label{lem:lim:pairs}
Define rational pairs $(u_j,v_j)$ by $(0,\frac14)$ for the constant $0$,
$(\frac14,\frac14)$ for the constant $1$, and, for $A_j=A_l\pm A_m$ and
$A_j=A_lA_m$ respectively,
\[
 u_j=\frac{u_lv_m\pm u_mv_l}4,\quad v_j=\frac{v_lv_m}4,\qquad
 u_j=\frac{u_lu_m}4,\quad v_j=\frac{v_lv_m}4 .
\]
Then $v_j>0$, $u_j=A_jv_j$, and $\abs{u_j},\abs{v_j}\le\frac14$ for every $j$.
The number $t=(2u_s-v_s)/4=v_s(2A-1)/4$ is nonzero, it is positive if and
only if $A>0$, and $\abs t\le\frac3{16}$.
\end{lemma}

\begin{proof}
By induction, $v_j$ is a product of positive numbers, and the ratio
$u_j/v_j$ equals $u_l/v_l\pm u_m/v_m$ or $(u_l/v_l)(u_m/v_m)$, which is
$A_j$. If all
earlier coordinates lie in $[-\frac14,\frac14]$, the new numerators have
absolute value at most $2\cdot\frac1{16}\cdot\frac14=\frac1{32}$ or
$\frac1{64}$. Since $A$ is an integer, $2A-1\ne0$, and $v_s>0$.
\end{proof}

List the coordinates $u_1,v_1,\ldots,u_s,v_s,t$ as $x_1,\ldots,x_N$, $N=2s+1$.
Each is given by a rule $x_i=p_i(x_1,\ldots,x_{i-1})$, where $p_i$ is a
constant, one of the bilinear formulas of \cref{lem:lim:pairs}, or the
affine formula for $t$. On the box $\mathcal B=[-\frac14,\frac14]^N$ these
rules satisfy
\begin{equation}\label{eq:lim:gate-bounds}
 \abs{p_i}\le\tfrac14,\qquad\norm{\nabla p_i}_1\le1,\qquad
 \norm{\nabla^2p_i}_2\le1 .
\end{equation}
Indeed, each bilinear rule has at most four gradient entries of absolute
value at most $\frac1{16}$, or fewer and larger entries when $l=m$, with
$\norm{\nabla p_i}_1\le\frac14$; its Hessian has at most two nonzero entries
in each row, of absolute value at most $\frac14$, or a single entry at most
$\frac12$ when $l=m$; and the rule for $t$ is affine with
$\norm{\nabla p_N}_1=\frac34$. The exact assignment $x^*$, defined by
$x_i^*=p_i(x_{<i}^*)$, lies in $\mathcal B$ by \cref{lem:lim:pairs}, and
$x_N^*=t^*$ has the sign of $A-\frac12$.

\begin{lemma}[Weighted gate objective]\label{lem:lim:gates}
Let $G(x)=\sum_{i=1}^N64^{N-i}\bigl(x_i-p_i(x)\bigr)^2$. Then
$\nabla^2G\succeq I$ on $\mathcal B$, $x^*$ is the unique zero of $G$ in
$\mathcal B$, and $G(x)\ge\frac12\norm{x-x^*}^2$ on $\mathcal B$.
\end{lemma}

\begin{proof}
Write $a_i=64^{N-i}$, $\rho_i=x_i-p_i(x)$, $D=\diag(\sqrt{a_i})$, and let $L$
be the strictly lower-triangular matrix with $L_{ij}=\partial_jp_i$. For a
direction $h$,
\[
 h^{\mathsf T}\nabla^2Gh=2\norm{D(I-L)h}^2-2\sum_ia_i\rho_i\,h^{\mathsf T}\nabla^2p_ih .
\]
Put $\widehat E=DLD^{-1}$, so $\widehat E_{ij}=8^{j-i}L_{ij}$ for $j<i$. By
\eqref{eq:lim:gate-bounds} every row of $\widehat E$ has absolute sum at
most $\frac18$, and every column has absolute sum at most
$\sum_{l\ge1}8^{-l}=\frac17$. Hence
$\norm{\widehat E}_2\le(\frac18\cdot\frac17)^{1/2}<\frac14$ and
\[
 2\norm{D(I-L)h}^2=2\norm{(I-\widehat E)Dh}^2\ge\tfrac98\norm{Dh}^2 .
\]
On $\mathcal B$, $\abs{\rho_i}\le\frac12$, and $p_i$ depends only on
$x_{<i}$, so $\abs{h^{\mathsf T}\nabla^2p_ih}\le\norm{h_{<i}}^2$. Therefore the
second term is at least
\[
 -\sum_ia_i\norm{h_{<i}}^2=-\sum_jh_j^2\sum_{i>j}a_i\ge-\frac1{63}\sum_ja_jh_j^2 .
\]
Altogether $h^{\mathsf T}\nabla^2Gh\ge(\frac98-\frac1{63})\norm{Dh}^2\ge\norm h^2$,
since every $a_j\ge1$. Because the rules are triangular, $G(x)=0$ exactly
when $x=x^*$. All residuals vanish at $x^*$, so $\nabla G(x^*)=0$, and
Taylor's theorem on the segment from $x^*$ to $x$ gives the last claim.
\end{proof}

\begin{lemma}[Two sign tests]\label{lem:lim:signs}
Let $\varepsilon=\frac1{32}$ and, on $\mathcal B\times[0,1]$, let
\[
 H_+(x,a)=G(x)+a^2-\varepsilon ax_N,\qquad H_-(x,a)=G(x)+a^2+\varepsilon ax_N .
\]
Each is strongly convex with modulus $1-\varepsilon$ and has a unique
minimizer $(x^\pm,a_\pm)$. Moreover $a_+>0$ if and only if $t^*>0$, and
$a_->0$ if and only if $t^*<0$. Exactly one of $a_+,a_-$ is positive.
\end{lemma}

\begin{proof}
The Hessian of $G(x)+a^2$ is at least $I$, and the bilinear term has Hessian
of operator norm $\varepsilon$; this gives the modulus and uniqueness on the
compact convex box. On the face $a=0$ both functions equal $G$, whose unique
minimizer on $\mathcal B$ is $x^*$, with $\nabla G(x^*)=0$. At $(x^*,0)$ the
$x$-gradient of $H_\pm$ is $\nabla G(x^*)=0$ and the $a$-derivative is
$\mp\varepsilon t^*$. If this derivative is nonnegative, $(x^*,0)$
satisfies the first-order conditions of the convex problem and is the
minimizer, so $a_\pm=0$. If it is negative, increasing $a$ decreases the
objective, so $(x^*,0)$ is not the minimizer; and no other point of the face
$a=0$ is, because a minimizer on that face would minimize $G$. Hence
$a_\pm>0$. Since $t^*\ne0$, exactly one case gives a positive value.
\end{proof}

\begin{lemma}[A convex amplifier]\label{lem:lim:amplifier}
Let $\eta=\frac1{192}$ and, on $Y=(\mathcal B\times[0,1])^2\times[0,1]$, let
\[
 \Phi(x^+,a_+,x^-,a_-,y)=H_+(x^+,a_+)+H_-(x^-,a_-)
 +\eta\bigl[a_+^2(y-1)^2+a_-^2y^2\bigr].
\]
Then $\Phi$ is convex on $Y$ and has a unique minimizer, whose
$y$-coordinate is $1$ if $t^*>0$ and $0$ if $t^*<0$.
\end{lemma}

\begin{proof}
For constants $c$ and a direction $(p,q)$ in the variables $(a,y)$, direct
differentiation gives
\[
 (p,q)^{\mathsf T}\nabla^2\bigl[a^2(y-c)^2\bigr](p,q)
 =2\bigl(aq+2(y-c)p\bigr)^2-6(y-c)^2p^2 .
\]
With $c=1$ for $a_+$ and $c=0$ for $a_-$, and $\abs{y-c}\le1$, the
amplifier contributes at least $-6\eta(p_+^2+p_-^2)$, where $p_\pm$ are the
components of the direction along $a_\pm$. The Hessian of $H_++H_-$ is at
least $(1-\varepsilon)I$ on all base coordinates. So the quadratic form of
$\nabla^2\Phi$ is at least $(1-\varepsilon-6\eta)\norm{p_{\mathrm{base}}}^2
=\frac{15}{16}\norm{p_{\mathrm{base}}}^2\ge0$, and $\Phi$ is convex.

The amplifier is nonnegative, so $\Phi\ge\min H_++\min H_-$. If $t^*>0$,
\cref{lem:lim:signs} gives $a_+>0=a_-$, and choosing the two base
minimizers and $y=1$ makes the amplifier vanish; so this lower bound is the
minimum. Every minimizer must then attain both base minima, which fixes all
base coordinates, and must make the amplifier vanish, which forces $y=1$
because $a_+>0$. The case $t^*<0$ is symmetric and forces $y=0$.
\end{proof}

\begin{proof}[Proof of \cref{thm:lim:posslp}]
The map writes $\Phi$ of \cref{lem:lim:amplifier}: two copies of the
$N$ local rules, the weights $64^{N-i}$ of $O(N)$ bits, and constant-size
remaining coefficients. Every residual $x_i-p_i(x)$ has at most three
monomials, so $\Phi$ is an explicit rational quartic with $O(N)$ monomials,
computed in time polynomial in the size of the program. By
\cref{lem:lim:amplifier,lem:lim:pairs}, $\Phi$ is convex on the rational box
$Y$, has a unique minimizer $\hat x$, and $\hat x_y=1$ exactly when $A>0$.

For $F_\gamma$, the core and residual terms are separate. The core term
$(v-\frac12)^2+\gamma v$ equals $(v-v_\gamma)^2$ plus a constant, so the
core optimizer is $v_\gamma=(1-\gamma)/2\in[\frac14,\frac34]$, the projected
growth is one, and $\partial_{vv}F_\gamma=2$. The residual objective $\Phi$
does not depend on $\gamma$, so the unique optimizer is
$(v_\gamma,\hat x)$ on every draw.

For (a), given a program, construct $\Phi$, compute the law, draw $\gamma$,
run the algorithm, and compute a point within distance $1/4$ of
$(v_\gamma,\hat x)$. Its $y$-coordinate is within $1/4$ of $\hat x_y$, so
comparing it with $\frac12$ decides whether $A>0$. The answer is correct on
every draw, and the expected running time is polynomial. For (b), apply the
deterministic algorithm to $\Phi$ directly. Square Root Sum is decidable in
polynomial time with a $\PosSLP$ oracle~\citep{allender2009-on-the-complexity-of-numerical}, so
$\PosSLP\in\mathrm P$ implies that Square Root Sum is in $\mathrm P$.
\end{proof}

\subsection{Proof of \texorpdfstring{\cref{prop:lim:flow}}{the flow-boundary proposition}}\label{app:lim:flow}

The atoms of $U_{1,M}$ are $-1+2j/(M-1)$, $0\le j<M$. Such an atom is at
least $\frac12$ exactly when $j\ge3(M-1)/4$. For a power of two $M\ge4$,
$3M/4$ is an integer and $3M/4-1<3(M-1)/4\le3M/4$, so the condition is
$j\ge3M/4$, which holds for $M/4$ indices. Each coordinate therefore has
probability $\frac14$, and by independence $\Pr(\mathcal E)=\frac1{16}$. The
second derivatives of $F_\gamma$ in the core are $\partial_{v_iv_i}=1$, so
$L=1$ is valid.

\emph{Part \ref{it:lim:flow-growth}.} Minimizing $v_1z_1+v_2z_2$ over the
two labels gives $\min\{v_1,v_2\}$, so $V_\gamma$ has the stated form. On
$\mathcal E$ each $\gamma_i>0$, and $v_i\ge v_i^2$ on $[0,1]$, so
$\gamma^{\mathsf T}v\ge\min_i\gamma_i\norm v^2$; with $\min\{v_1,v_2\}\ge0$
this gives $V_\gamma(v)\ge(\frac12+\min_i\gamma_i)\norm v^2\ge\norm v^2$.
Since $V_\gamma(0)=0$, the origin is the unique optimal core, and both
labels have cost zero there.

\emph{Part \ref{it:lim:flow-box}.} At $(a,0)$ with $a>0$ the label $(1,0)$
costs $a$ and $(0,1)$ costs $0$; at $(0,b)$ with $b>0$ the opposite holds.
The noise term is common to both labels and cancels from these comparisons.
So each label is strictly worse than the other at some point of the box,
and no valid certificate can assert that one label is optimal throughout.

\emph{Part \ref{it:lim:flow-search}.} The cell $[0,h]^2$ has the corner $0$
with value $V_\gamma(0)=\min V_\gamma$, and the incumbent $U_j$ is a
feasible value, hence at least $\min V_\gamma$. So the lower bound of the
cell, at most $V_\gamma(0)-h^2/4$, does not exceed $U_j$, and the cell is
retained whenever its parent is. At level $0$ the whole core box is
retained, and induction gives the claim at every level.

\emph{The chain.} With $m-1$ further stages of two zero-cost parallel arcs,
each feasible flow chooses one arc per stage, so there are $2^m$ flows. The
conditional value is unchanged, so parts
\ref{it:lim:flow-growth}--\ref{it:lim:flow-search} hold, and on
$\mathcal E$ every flow is optimal at the origin. In every box
$[0,a]\times[0,b]$ the optimal flows at $(a,0)$ and at $(0,b)$ use
different first-stage arcs, so no single flow is certified on the retained
hull at any level. A strategy that enumerates all flows when this
certificate fails does at least $2^m$ work on $\mathcal E$, and its
expected work is at least $2^m/16$. The network has $2m$ unit-capacity arcs
and $m+1$ nodes, so its encoding has length $O(m\log m)$.\qed
